% Subsection 1: Design Goals and Abstraction Level
\subsection{Design Goals and Abstraction Level}
% we want to mirror arch intent, but at what level
% what atalla-sim wants to accomplish and at what fidelity
% purpose: pre-silicon arch exploration
% questions it wants to answer: cycle counts, queue pressure, overlapping data movement and compute, utilization, sensitivity to uarch choices
% define abstraction level: clock-level simulation for modeled paths, full RTL reproduction
% modeling philosophy: detail behaviors that affect perf and correctness, abstract away things that dont
% intended boundaries: no full OS/runtime model, no phys design, some harness-driven policies instead of full prod scheduler
Atalla-Sim is a pre-silicon architectural simulator for exploring the Atalla design space before RTL is complete. It collects cycle counts, queue pressure, utilization, FLOP counts, and bytes transmitted so that design choices can be evaluated in terms of execution time, starvation behavior, overlap of data movement with computation, and likely bottlenecks. The simulator is therefore detailed where behavior changes architectural conclusions and deliberately approximates where it does not.

This trade-off defines the simulator's abstraction level. Behaviors that materially affect correctness, cycle counts, or contention are modeled with clock-level fidelity, including instruction issue, queue transitions, memory latency, inter-unit handshakes, and architectural writeback. Details such as wire timing, power-gating, and physical floorplanning are omitted. Besides that, some subsystems use calibrated abstractions rather than full reproduction. For example, DRAM timing can be parameterized from higher-fidelity tools such as Ramulator\cite{ramulator}. Atalla-Sim is therefore an architectural exploration and hardware-software co-design tool, not a full verification environment.

% Subsection 2: Runtime Platform and Core Classes
\subsection{Runtime Platform and Core Classes}
% what are major sim objects, how do they fit together
% clocked object
% top-level platform
% vector core as center of issue, execution, and writeback
% execution side with vdp, VectorRegisterFile, gsau, vls, sysarr
% mem side with DRAM, scpad
% bridges

The Atalla-Sim runtime platform is organized as a network of modeled hardware objects that advance through per-cycle \texttt{tick} invocations. Every timed component inherits from the \texttt{Clocked} base class, which provides a uniform interface for cycle advancement and keeps all state updates inside the simulation loop. The platform is queue-centric: instructions, operands, and memory requests move between components through architectural queues, so the runtime organization directly mirrors the abstracted hardware pipeline shown in \autoref{fig:runtimeplatfordataflowchart}.

The \texttt{VectorCore} sits at the center of the platform and coordinates instruction dispatch, operand collection, execution, and writeback. On the compute side, \texttt{VectorDatapath} handles lane-parallel operations, \texttt{VectorRegisterFile} holds architectural vector state, and \texttt{GSAU} with \texttt{GSAUSystolicArrayBridge} stages traffic into \texttt{SystolicArray} and returns completed rows to shared writeback. On the memory side, \texttt{Scratchpad} provides software-managed on-chip storage, \texttt{DRAM} models off-chip latency, \texttt{VLSFrontendBridge} connects the vector load/store path to the scratchpad frontends, and \texttt{Backend} models DRAM-facing transaction flow. Workload harnesses such as \texttt{TiledAtallaCosim} and \texttt{MNReuseBlockedAtallaCosim} sit above this reusable platform to stage data, inject tiled work, and verify
\begin{landscape}
\begin{figure}
    \centering
    \includegraphics[width=\linewidth,height=.95\textheight,keepaspectratio]{figures/runtime_platform_data_flowchart_drawio.png}
    \caption{Runtime Platform Data Flowchart}
    \label{fig:runtimeplatfordataflowchart}
\end{figure}
\end{landscape}

completion without hard-wiring the workload policy into the machine model.

% Subsection 3: Packetized Execution Model
\subsection{Packetized Execution Model}
% how does sim execute vliw behavior
% packet in sim terms
% instrs collected in a build packet before entering packet queue
% per-unit struct inside packet: gsau, vls, vdp entries
% slot constraints and width limits for each unit type
% how packets are consider cycle by cycle, split into per-unit issue paths
% structural hazards/backpressure: packets only partially issued
% packet complete and removed from queue
% WB completion after execution
The \texttt{VectorCore} scheduler uses a packetized instruction model that captures VLIW-style issue timing without reproducing the full production scheduler. Instructions accumulate in \texttt{\_build\_packet}, which holds separate per-unit lists for the GSAU path, the VLSU path, and the vector datapath. The current slot limits are one GSAU instruction, up to four VLSU instructions total, and up to two datapath instructions. A packet is flushed into \texttt{vliw\_q} when the next instruction would exceed one of those limits, and any partially built packet can also be flushed before issue.

Each cycle, \texttt{VectorCore} processes the head packet in \texttt{vliw\_q} and attempts to issue its instructions to the GSAU, VLSU, and datapath fronts in parallel. If one path is back-pressured, that instruction remains in the packet while ready instructions for other paths can still issue. Once a datapath instruction enters \texttt{VectorDatapath}'s pending queue, later hazards are modeled inside the datapath rather than in \texttt{vliw\_q}. The completed results from the VLSUs, the GSAU return path, and the datapath converge in \texttt{WBBuffer}, and \texttt{VectorCore} commits at most one architectural writeback per cycle to \texttt{VectorRegisterFile}.

% Subsection 4: Memory Hierarchy and Data Movement
\subsection{Memory Hierarchy and Data Movement}
% how data moves through sim system, latency and contention
% full modeled dataflow dram -> scpad -> vc -> sysarr
% dram + backend, burst based, tx queuing
% scpad org
% vls mem ops -> req/rsp
% data is staged in VectorRegisterFile
% VectorRegisterFile -> gsau -> sysarr
% return + accumulation
% contention points: q depth, bank conflict, frontend saturation, backend serialization, unit acceptance limits
Data movement in Atalla-Sim follows the modeled path \texttt{DRAM} $\rightarrow$ \texttt{Scratchpad} $\rightarrow$ \texttt{VectorRegisterFile} $\rightarrow$ compute fabric $\rightarrow$ writeback. Each stage contributes latency and potential contention, which determines how well data movement overlaps with compute.

Off-chip memory is modeled by \texttt{DRAM} and \texttt{Backend}. The backend is burst-oriented, with queued and in-flight transactions, so refill pressure is limited by channel serialization. Returned data enters \texttt{Scratchpad}, whose multiple frontends, crossbars, and \texttt{SRAMBanks} expose both frontend saturation and bank-conflict effects.

From there, data enter the VLS path through \texttt{VLSFrontendBridge} and are handled by the \texttt{VLSU} through issue, request, response, destination tracking, and writeback into \texttt{WBBuffer}. Loads become architecturally visible only when \texttt{VectorCore} commits them into \texttt{VectorRegisterFile}. GSAU instructions then move operands from \texttt{VectorRegisterFile} through \texttt{GSAU} and \texttt{GSAUSystolicArrayBridge} into \texttt{SystolicArray}, and completed rows return through the same shared writeback path. Along this route, the simulator records cycle counts, queue occupancies, bandwidth, and utilization so that throughput losses can be localized.

% Subsection 5: Harness Design and Workload Injection
\subsection{Harness Design and Workload Injection}
% talk about overall test harness design, not the specific workload we are running (1024x1024 gemm)
% how workload is driven on top of sim platform
% harness assembles platform, initializes mem, inject workloads, collects measurements
% reusable platfotm-building helps vs workload-specific harness logic
% how harness drives the sim, checks for completion
% what metrics are we getting
Workloads in Atalla-Sim are driven by a harness layer above the runtime platform. The harness assembles the simulated system, initializes memory, injects work, advances the machine, and collects architectural metrics. \texttt{TiledAtallaCosim} and \texttt{MNReuseBlockedAtallaCosim} are representative examples in the current Atalla workflow.

Platform assembly is handled through reusable helpers such as \texttt{build\_atalla\_platform(\ldots)}, which instantiate and connect the common runtime objects. This keeps the machine model reusable while leaving workload policy in the harness. During execution, the harness advances the vector core, bridges, backends, and scratchpad in dependency order while monitoring queue and pipeline state to decide when to inject new work and when the machine has drained. Because workload latency is variable, completion is detected by idleness in the relevant queues and shared writeback structures rather than by a fixed cycle budget.

At the end of a run, the harness extracts the architecture-facing counters used throughout the evaluation, including total cycles, DRAM-side and scratchpad/vector-core bytes transmitted, queue-depth summaries, bandwidth estimates, functional-unit utilization, and other accumulated platform counters.